\CWHeader{Лабораторная работа \textnumero 8}

\CWProblem{
Бычкам дают пищевые добавки, чтобы ускорить их рост. Каждая добавка содержит некоторые из $N$ действующих веществ. Соотношения количеств веществ в добавках могут отличаться.

Воздействие добавки определяется как $c_1 a_1 + c_2 a_2 + \dots + c_N a_N$, где $a_i$~--- количество $i$-го вещества в добавке, $c_i$~--- неизвестный коэффициент, связанный с веществом и не зависящий от добавки. Чтобы найти неизвестные коэффициенты $c_i$, Биолог может измерить воздействие любой добавки, использовав один её мешок.

Известна цена мешка каждой из $M$ ($M \le N$) различных добавок. Нужно помочь Биологу подобрать самый дешёвый набор добавок, позволяющий найти коэффициенты $c_i$. Возможно, соотношения веществ в добавках таковы, что определить коэффициенты нельзя.

\textbf{Формат ввода}

В первой строке записаны целые числа $M$ и $N$. В каждой из следующих $M$ строк записаны $N$ чисел, задающих соотношение количеств веществ в добавке, а за ними~--- цена мешка добавки. Порядок веществ во всех описаниях добавок один и тот же. Все числа~--- неотрицательные целые, не превосходящие 50.

\textbf{Формат вывода}

Выведите $-1$, если определить коэффициенты невозможно. Иначе выведите набор добавок в виде их номеров (добавки нумеруются по порядку во входных данных). Если вариантов несколько, выведите любой из них.
}
\pagebreak